../cictikz/src/cictikz/data/tex/cictikz_preamble.tex

% Thermometer coded resistor DAC: every element is the same R out of -V_REF,
% and each one is either steered into the summing node or into ground. Three
% equal elements rather than a binary weighted set is the whole point — the
% code from dac_thermo_states can only ever add or remove one of them, so
% there is no transition where one element leaves as another arrives.

\begin{circuitikz}[american, thick, transform shape, circuit ee IEC]

  ../cictikz/src/cictikz/data/tex/cictikz_lib.tex
  ../cictikz/src/cictikz/data/tex/rdac_lib.tex

  \def\ybot{0.4}
  \def\ytop{4.8}

  % ---- Three equal elements on the -V_REF rail ----
  \foreach \x in {0, 2.2, 4.4} {
    \draw (\x,\ybot) \vresistor{$R$} -- ++(0,0.7);
  }
  \draw (0,\ybot) -- (4.4,\ybot);
  \draw (0,\ybot) -- ++(-1.2,0) coordinate (vrefport);
  \rdacPort{(vrefport)}{north}{$-V_{REF}$}

  % ---- Each element steers into the summing rail or into ground ----
  \foreach \x in {0, 2.2, 4.4} {
    \rdacSpdt{(\x,{\ybot+2.3})}{\ytop}{}{}
  }

  % ---- The rail runs into the amplifier ----
  \draw (0.55,\ytop) -- (8.8,\ytop);
  \fill (2.75,\ytop) circle (0.07);
  \fill (4.95,\ytop) circle (0.07);
  \rdacAmp{OA}{(8.8,4.2)}
  \draw (8.8,\ytop) -- (OAinn);

  % ---- Feedback resistor over the top ----
  \draw (8.0,\ytop) -- (8.0,7.0) -- ++(0.7,0) \hresistor{$R_F$} -- ++(0.7,0)
        -- (12.0,7.0) -- (12.0,4.2) -- (OAout);
  \fill (8.0,\ytop) circle (0.07);
  \draw (OAout) -- ++(0.9,0);

  % ---- The + input sits on ground ----
  \draw (OAinp) -- ++(-1.2,0) coordinate (og);
  \draw (og) -- ++(0,-0.3) \vground;

\end{circuitikz}
\end{document}
